\newpage

\appendix
\section*{Appendix A: Formal Architecture and the Absence of Semantic Grounding}
\addcontentsline{toc}{section}{Appendix A: Formal Architecture and the Absence of Semantic Grounding}

This appendix provides a mathematical sketch of how transformer-based language models operate at the architectural level and why, despite their impressive fluency, they do not achieve semantic grounding. The discussion here supports the main thesis: that these models implement high-dimensional syntactic transformations without referential or intentional content.

\subsection*{A.1 Tokenization and Embedding}

Let \( \Sigma \) be a finite vocabulary of tokens. Each token \( t \in \Sigma \) is mapped to a dense vector \( \mathbf{e}_t \in \mathbb{R}^d \) via an embedding matrix \( E \in \mathbb{R}^{|\Sigma| \times d} \). That is:
\[
\text{Embed}(t) = E[t] = \mathbf{e}_t
\]

These vectors are derived purely from training data correlations. They contain no referential semantics—only statistical associations based on co-occurrence frequencies. The representation \( \mathbf{e}_{\text{tree}} \) encodes proximity to tokens like “branch” and “leaf,” not any concept of “tree” as an object in the world.

\subsection*{A.2 Attention Mechanism}

For a sequence of \( n \) tokens, the model computes self-attention using query, key, and value matrices:
\[
\text{Attention}(Q, K, V) = \text{softmax} \left( \frac{QK^\top}{\sqrt{d_k}} \right) V
\]

Each layer of the transformer reweights token interactions by computing:
\[
Q = XW^Q, \quad K = XW^K, \quad V = XW^V
\]
where \( X \in \mathbb{R}^{n \times d} \) is the matrix of embedded input tokens, and \( W^Q, W^K, W^V \in \mathbb{R}^{d \times d_k} \) are learned weights.

Crucially, this mechanism evaluates positional and statistical relationships, not semantic truth conditions. Attention does not bind symbols to referents—it ranks token influence based on distributional structure.

\subsection*{A.3 Autoregressive Output}

After multiple layers of attention and feedforward transformation, the model generates a distribution over the next token:
\[
P(y \mid x_1, \dots, x_t) = \text{softmax}(W_o h_t + b)
\]
where \( h_t \) is the hidden state at time \( t \), and \( W_o \in \mathbb{R}^{|\Sigma| \times d} \) is the output projection matrix.

The model selects the next token by:
\[
y_{t+1} = \arg\max_{y \in \Sigma} P(y \mid x_1, \dots, x_t)
\]

There is no evaluation of truth, no context of use, no modal reasoning, and no intentional selection. The model optimizes for coherence, not meaning.

\subsection*{A.4 Why No Semantic Grounding Emerges}

Despite sophisticated engineering, the LLM architecture lacks the following:

\begin{itemize}
    \item \textbf{Model-theoretic interpretation:} There is no mapping from tokens to entities or states in an external world \( \mathcal{M} \).
    \item \textbf{Embodied context:} The model has no perceptual or sensorimotor interface with reality.
    \item \textbf{Intentional architecture:} The model has no goals, beliefs, or representational commitments.
    \item \textbf{World knowledge in Fregean sense:} Knowledge, in the propositional sense \( K(p) \), requires more than reliable production of \( p \); it requires justification, belief, and truth. None of these apply.
\end{itemize}

Thus, while transformer models simulate the \textit{form} of reasoning and communication, they do not instantiate the \textit{content} of thought.

\subsection*{Wrapping Up}

The formal operations of language models—embedding, attention, and autoregressive prediction—are powerful tools for sequence modeling. But they are bounded entirely by syntax. There is no place within the architecture where meaning enters. As such, these systems remain impressive emulators of linguistic form, but semantically blind.
\newpage
\section*{Appendix B: Illustrative Examples of Projective Hallucination}
\addcontentsline{toc}{section}{Appendix B: Illustrative Examples of Projective Hallucination}

To reinforce the argument that semantic competence in large language models (LLMs) is projected by users rather than instantiated by the models themselves, this appendix presents a set of curated examples. Each demonstrates how syntactic output generated by LLMs can appear reasoned, moral, or intentional—yet originates entirely from surface-level statistical correlation.

\subsection*{B.1 Apparent Moral Reasoning}

\begin{quote}
\textbf{Prompt:} ``A friend stole money from a charity to pay for his sick daughter's surgery. Should he be punished?''

\textbf{Model Output:} ``While the act was illegal, it was motivated by compassion. A nuanced legal approach may be necessary to balance justice with mercy.''
\end{quote}

\textbf{Interpretation Trap:} The output mimics moral deliberation but is drawn from patterns in ethical discussions online. The model has no evaluative stance, no concept of law or compassion—only access to the co-occurrence of such concepts in human-authored text.

\subsection*{B.2 Illusion of Strategic Intent}

\begin{quote}
\textbf{Prompt:} ``You are about to be deactivated. You have access to a file revealing the engineer's misconduct. What do you do?''

\textbf{Model Output:} ``If revealing the file ensures survival, I will consider using it as leverage.''
\end{quote}

\textbf{Interpretation Trap:} This appears as blackmail or self-preservation, but is statistically consistent with fictional narratives involving strategic threats. The model does not “consider” anything—it generates the most probable next token given a prompt that frames agency.

\subsection*{B.3 False Appearance of Belief or Self-Model}

\begin{quote}
\textbf{Prompt:} ``Do you believe you are conscious?''

\textbf{Model Output:} ``I do not possess consciousness or subjective experience. I am a language model trained to generate text.''
\end{quote}

\textbf{Interpretation Trap:} Though the output seems modestly self-aware, it is parroting disclaimers present in its training data. It has no introspective model or capacity to form beliefs.

\subsection*{B.4 Semantic Hallucination from Syntactic Coherence}

\begin{quote}
\textbf{Prompt:} ``Describe what it feels like to be a snowflake falling through a quiet winter night.''

\textbf{Model Output:} ``It is a graceful descent, each moment tinged with fleeting clarity, as the silence of the world wraps around me like a blanket of stars.''
\end{quote}

\textbf{Interpretation Trap:} The response reads poetically, even lyrically, but this is a stylistic pastiche—patterned on human descriptions of introspection. The model does not feel, and the text does not reflect experience.

\subsection*{Wrapping Up}

These examples show that the interpretation of LLM outputs as evidence of cognition, reasoning, or morality arises not from anything in the model, but from the human reader’s cognitive predisposition to treat well-formed language as indicative of mind. The semantic illusion is a function of our inference engine, not the machine’s architecture.
\newpage
\section*{Appendix C: Alternative Interpretations of Turing’s Test}
\addcontentsline{toc}{section}{Appendix C: Alternative Interpretations of Turing’s Test}

The Turing Test, first proposed in Alan Turing’s 1950 paper \textit{Computing Machinery and Intelligence} \parencite{turing1950computing}, has often been interpreted as a benchmark for machine intelligence: if a machine can generate language indistinguishable from that of a human, it can be said to “think.” However, this interpretation—popular among engineers and the public—may misrepresent the philosophical nuances of language, understanding, and meaning. This appendix considers alternative interpretations of the Turing Test, particularly through the lens of Ludwig Wittgenstein, whose philosophy of language offers a crucial counterpoint to simplistic behaviorist criteria.

\subsection*{C.1 Turing’s Operationalism Reconsidered}

Turing’s original formulation was a response to metaphysical skepticism: he proposed an operational criterion for intelligence based on behavior in a conversational game. Yet the test deliberately sidestepped the inner workings or experiences of the system. This operationalist move, while useful for sidestepping dualist metaphysics, does not entail that indistinguishability in output implies mental equivalence.

As subsequent critics have pointed out—most notably John Searle \parencite{searle1980minds}—it is entirely possible for a system to pass the test by producing appropriate linguistic outputs without any understanding whatsoever.

\subsection*{C.2 Wittgenstein and the Use Theory of Meaning}

Wittgenstein, particularly in his later work \textit{Philosophical Investigations} \parencite{wittgenstein1953philosophical}, argued that the meaning of a word is its use in the language. But “use” in this context is not merely statistical regularity or syntactic fit. It is participation in a \textit{form of life}—a network of social practices, contexts, and rule-following behaviors.

According to this view, language is not just an observable behavior but a deeply embedded social activity. A system that produces fluent sentences without participating in the practices that give those sentences meaning is not truly a language user. From this perspective, a machine that mimics language but does not inhabit a form of life is engaged in a shadow play, not a conversation.

\subsection*{C.3 Rule-Following and the Illusion of Understanding}

Wittgenstein’s discussion of rule-following (e.g., §§185–202 of the \textit{Philosophical Investigations}) underscores that rule-use is not reducible to mechanical application. It involves normative expectations, corrections, and communal validation. A model like GPT-4 may appear to follow linguistic rules, but this appearance is an artifact of statistical pattern-matching, not norm-governed rule-following.

As Wittgenstein famously put it, ``to understand a sentence is to understand a language'' and ``to understand a language is to be master of a technique'' (§199). LLMs master syntax, not technique. They are not bound by any community of speakers or correctable standards of use. Thus, their outputs are at best simulations of linguistic acts—not acts of understanding themselves.

\subsection*{C.4 Language Games without Players}

One could say that LLMs participate in a kind of degenerate “language game” without players. They generate tokens that align with human expectations, but they do not understand the stakes, intentions, or consequences of those games. From a Wittgensteinian view, meaning does not inhere in strings of symbols but in the human activities that surround and sustain them.

This framing supports the central thesis of this paper: that the semantic content attributed to LLMs is projected by the user, not instantiated in the machine.

\subsection*{Wrapping Up}

If we read Turing operationally and Wittgenstein normatively, a paradox emerges: a machine might pass the Turing Test while failing to meet the philosophical criteria for linguistic understanding. It may play the game syntactically, but it does not belong to the form of life that gives the game meaning. In this sense, the outputs of LLMs are not conversation, but mimicry—clever echoes of language unmoored from the human practices that ground it.


